% GENERATED FILE. Edit canonical lessons, then run npm run generate:books.
% canonical-source-hash: fnv1a64:02a07e61d73f6996
% canonical-lessons: KA-C410-higagi, KA-C410-adagyu, KA-C410-nimitta, KA-C410-melagi, KA-C410-matravallade

\chapter{Linking Ideas}
\label{ch:ka-linking-ideas}
\hlchaptermodality{410}

\begin{chapteropening}
\textbf{By the end of this chapter:} \emph{I can say thus, however, on account of, above all and not only that.}

Complete the last lesson of chapter 410: I can say thus, however, on account of, above all and not only that.
\end{chapteropening}

\section[\texorpdfstring{\kn{ಹೀಗಾಗಿ} (hīgāgi)}{ಹೀಗಾಗಿ (hīgāgi)}]{\kn{ಹೀಗಾಗಿ} (hīgāgi) — thus; for this reason}

\label{lesson:KA-C410-higagi}



\begin{quote}
Before the new one: say the Kannada for to present, then the Kannada for a note.
\end{quote}

\subsection*{You'll want to know: \kn{ಹೀಗಾಗಿ}}
\textbf{\kn{ಹೀಗಾಗಿ}} — \emph{hīgāgi} — \textquotedblleft{}thus; for this reason\textquotedblright{}.

\begin{quote}
\kn{ಬಸ್} \kn{ತಡವಾಯಿತು}, \kn{ಹೀಗಾಗಿ} \kn{ನಾನು} \kn{ತಡವಾಗಿ} \kn{ಬಂದೆ}.
\end{quote}

\begin{quote}
\emph{bas taḍavāyitu, hīgāgi nānu taḍavāgi bande.}
\end{quote}

\begin{quote}
The bus was late, so I came late.
\end{quote}

\kn{ಹೀಗೆ} (like this) + \kn{ಆಗಿ} (being).

\begin{cousinweb}
\textbf{\kn{ಹೀಗಾಗಿ}} (\emph{hīgāgi}) is native Dravidian, a phrase grown into one word: \emph{hīge}, \textquotedblleft{}in this way\textquotedblright{} (DEDR 410), and \emph{āgi}, \textquotedblleft{}having become\textquotedblright{} (DEDR 333).
\end{cousinweb}

\begin{grammarlens}[title={What you've built}]
One more word for giving reasons and opinions.
\end{grammarlens}

\subsection*{Your turn}
\begin{itemize}
\raggedright
  \item \emph{Say it:} \emph{hīgāgi}
  \item \emph{Say it:} \emph{hīgāgi}, once more
  \item \emph{Say it:} \emph{bas taḍavāyitu, hīgāgi nānu taḍavāgi bande.} — The bus was late, so I came late.
  \item \emph{Recall:} say the Kannada for to present, then the Kannada for a note, then say \emph{hīgāgi} again
\end{itemize}

\begin{tcolorbox}[breakable,colback=teal!4,colframe=teal!35!black,title={Before you move on}]
What does \kn{ಹೀಗಾಗಿ} mean? (\textquotedblleft{}Thus; for this reason\textquotedblright{}.) Say it once more.
\end{tcolorbox}
\section[\texorpdfstring{\kn{ಆದಾಗ್ಯೂ} (ādāgyū)}{ಆದಾಗ್ಯೂ (ādāgyū)}]{\kn{ಆದಾಗ್ಯೂ} (ādāgyū) — however; nevertheless}

\label{lesson:KA-C410-adagyu}



\begin{quote}
Before the new one: say the Kannada for a note, then the Kannada for thus.
\end{quote}

\subsection*{You'll want to know: \kn{ಆದಾಗ್ಯೂ}}
\textbf{\kn{ಆದಾಗ್ಯೂ}} — \emph{ādāgyū} — \textquotedblleft{}however; nevertheless\textquotedblright{}.

\begin{quote}
\kn{ಅವನು} \kn{ಬಡವ}, \kn{ಆದಾಗ್ಯೂ} \kn{ಎಲ್ಲರಿಗೂ} \kn{ಸಹಾಯ} \kn{ಮಾಡುತ್ತಾನೆ}.
\end{quote}

\begin{quote}
\emph{avanu baḍava, ādāgyū ellarigū sahāya māḍuttāne.}
\end{quote}

\begin{quote}
He is poor; nevertheless he helps everyone.
\end{quote}

More formal than \kn{ಆದರೂ}.

\begin{cousinweb}
\textbf{\kn{ಆದಾಗ್ಯೂ}} (\emph{ādāgyū}) is native Dravidian, built on forms of \emph{āgu}, \textquotedblleft{}to become\textquotedblright{} (DEDR 333), with -\emph{ū}, \textquotedblleft{}even\textquotedblright{}: even so.
\end{cousinweb}

\begin{grammarlens}[title={What you've built}]
One more word for giving reasons and opinions.
\end{grammarlens}

\subsection*{Your turn}
\begin{itemize}
\raggedright
  \item \emph{Say it:} \emph{ādāgyū}
  \item \emph{Say it:} \emph{ādāgyū}, once more
  \item \emph{Say it:} \emph{avanu baḍava, ādāgyū ellarigū sahāya māḍuttāne.} — He is poor; nevertheless he helps everyone.
  \item \emph{Recall:} say the Kannada for a note, then the Kannada for thus, then say \emph{ādāgyū} again
\end{itemize}

\begin{tcolorbox}[breakable,colback=teal!4,colframe=teal!35!black,title={Before you move on}]
What does \kn{ಆದಾಗ್ಯೂ} mean? (\textquotedblleft{}However; nevertheless\textquotedblright{}.) Say it once more.
\end{tcolorbox}
\section[\texorpdfstring{\kn{ನಿಮಿತ್ತ} (nimitta)}{ನಿಮಿತ್ತ (nimitta)}]{\kn{ನಿಮಿತ್ತ} (nimitta) — on account of; an occasion}

\label{lesson:KA-C410-nimitta}



\begin{quote}
Before the new one: say the Kannada for thus, then the Kannada for however.
\end{quote}

\subsection*{You'll want to know: \kn{ನಿಮಿತ್ತ}}
\textbf{\kn{ನಿಮಿತ್ತ}} — \emph{nimitta} — \textquotedblleft{}on account of; an occasion\textquotedblright{}.

\begin{quote}
\kn{ಹಬ್ಬದ} \kn{ನಿಮಿತ್ತ} \kn{ಶಾಲೆಗೆ} \kn{ಎರಡು} \kn{ದಿನ} \kn{ರಜೆ} \kn{ಇದೆ}.
\end{quote}

\begin{quote}
\emph{habbada nimitta śālege eraḍu dina raje ide.}
\end{quote}

\begin{quote}
On account of the festival, the school has two days off.
\end{quote}

It follows a noun in -\kn{ದ} / -\kn{ಯ}: \kn{ಮದುವೆಯ} \kn{ನಿಮಿತ್ತ}, for the wedding. A formal word for a reason or an occasion.

\begin{cousinweb}
\textbf{\kn{ನಿಮಿತ್ತ}} (\emph{nimitta}) is borrowed from Sanskrit \emph{nimitta}, \textquotedblleft{}a cause, a reason\textquotedblright{} (CDIAL 7241); in Sanskrit, too, its case forms mean \textquotedblleft{}on account of\textquotedblright{} (Monier-Williams).
\end{cousinweb}

\begin{grammarlens}[title={What you've built}]
One more word for giving reasons and opinions.
\end{grammarlens}

\subsection*{Your turn}
\begin{itemize}
\raggedright
  \item \emph{Say it:} \emph{nimitta}
  \item \emph{Say it:} \emph{nimitta}, once more
  \item \emph{Say it:} \emph{habbada nimitta śālege eraḍu dina raje ide.} — On account of the festival, the school has two days off.
  \item \emph{Recall:} say the Kannada for thus, then the Kannada for however, then say \emph{nimitta} again
\end{itemize}

\begin{tcolorbox}[breakable,colback=teal!4,colframe=teal!35!black,title={Before you move on}]
What does \kn{ನಿಮಿತ್ತ} mean? (\textquotedblleft{}On account of; an occasion\textquotedblright{}.) Say it once more.
\end{tcolorbox}
\section[\texorpdfstring{\kn{ಮೇಲಾಗಿ} (mēlāgi)}{ಮೇಲಾಗಿ (mēlāgi)}]{\kn{ಮೇಲಾಗಿ} (mēlāgi) — above all; moreover}

\label{lesson:KA-C410-melagi}



\begin{quote}
Before the new one: say the Kannada for however, then the Kannada for on account of.
\end{quote}

\subsection*{You'll want to know: \kn{ಮೇಲಾಗಿ}}
\textbf{\kn{ಮೇಲಾಗಿ}} — \emph{mēlāgi} — \textquotedblleft{}above all; moreover\textquotedblright{}.

\begin{quote}
\kn{ಈ} \kn{ಮನೆ} \kn{ದೊಡ್ಡದು}, \kn{ಮೇಲಾಗಿ} \kn{ಬಾಡಿಗೆ} \kn{ಕಡಿಮೆ}.
\end{quote}

\begin{quote}
\emph{ī mane doḍḍadu, mēlāgi bāḍige kaḍime.}
\end{quote}

\begin{quote}
This house is big; moreover the rent is low.
\end{quote}

\kn{ಮೇಲೆ} (above) + \kn{ಆಗಿ}.

\begin{cousinweb}
\textbf{\kn{ಮೇಲಾಗಿ}} (\emph{mēlāgi}) is native Dravidian: \emph{mēl}, \textquotedblleft{}what is above\textquotedblright{} (DEDR 5086), with \emph{āgi}, \textquotedblleft{}being\textquotedblright{} (DEDR 333).
\end{cousinweb}

\begin{grammarlens}[title={What you've built}]
One more word for giving reasons and opinions.
\end{grammarlens}

\subsection*{Your turn}
\begin{itemize}
\raggedright
  \item \emph{Say it:} \emph{mēlāgi}
  \item \emph{Say it:} \emph{mēlāgi}, once more
  \item \emph{Say it:} \emph{ī mane doḍḍadu, mēlāgi bāḍige kaḍime.} — This house is big; moreover the rent is low.
  \item \emph{Recall:} say the Kannada for however, then the Kannada for on account of, then say \emph{mēlāgi} again
\end{itemize}

\begin{tcolorbox}[breakable,colback=teal!4,colframe=teal!35!black,title={Before you move on}]
What does \kn{ಮೇಲಾಗಿ} mean? (\textquotedblleft{}Above all; moreover\textquotedblright{}.) Say it once more.
\end{tcolorbox}
\section[\texorpdfstring{\kn{ಮಾತ್ರವಲ್ಲದೆ} (mātravallade)}{ಮಾತ್ರವಲ್ಲದೆ (mātravallade)}]{\kn{ಮಾತ್ರವಲ್ಲದೆ} (mātravallade) — not only that; besides}

\label{lesson:KA-C410-matravallade}



\begin{quote}
Before the new one: say the Kannada for on account of, then the Kannada for above all.
\end{quote}

\subsection*{You'll want to know: \kn{ಮಾತ್ರವಲ್ಲದೆ}}
\textbf{\kn{ಮಾತ್ರವಲ್ಲದೆ}} — \emph{mātravallade} — \textquotedblleft{}not only that; besides\textquotedblright{}.

\begin{quote}
\kn{ಅವಳು} \kn{ಹಾಡುತ್ತಾಳೆ} \kn{ಮಾತ್ರವಲ್ಲದೆ} \kn{ನೃತ್ಯವನ್ನೂ} \kn{ಮಾಡುತ್ತಾಳೆ}.
\end{quote}

\begin{quote}
\emph{avaḷu hāḍuttāḷe mātravallade nṛtyavannū māḍuttāḷe.}
\end{quote}

\begin{quote}
She not only sings but dances too.
\end{quote}

\kn{ಮಾತ್ರ} (only) + \kn{ಅಲ್ಲದೆ} (not being).

\begin{cousinweb}
\textbf{\kn{ಮಾತ್ರವಲ್ಲದೆ}} (\emph{mātravallade}) is built from \emph{mātra}, \textquotedblleft{}only\textquotedblright{}, from Sanskrit \emph{mātra}, \textquotedblleft{}measure\textquotedblright{} (Monier-Williams), joined to native \emph{allade}, \textquotedblleft{}not being\textquotedblright{} (DEDR 234).
\end{cousinweb}

\begin{grammarlens}[title={What you've built}]
One more word for giving reasons and opinions.
\end{grammarlens}

\subsection*{Your turn}
\begin{itemize}
\raggedright
  \item \emph{Say it:} \emph{mātravallade}
  \item \emph{Say it:} \emph{mātravallade}, once more
  \item \emph{Say it:} \emph{avaḷu hāḍuttāḷe mātravallade nṛtyavannū māḍuttāḷe.} — She not only sings but dances too.
  \item \emph{Recall:} say the Kannada for on account of, then the Kannada for above all, then say \emph{mātravallade} again
\end{itemize}

\begin{tcolorbox}[breakable,colback=teal!4,colframe=teal!35!black,title={Before you move on}]
What does \kn{ಮಾತ್ರವಲ್ಲದೆ} mean? (\textquotedblleft{}Not only that; besides\textquotedblright{}.) Say it once more.
\end{tcolorbox}
